% Rúbrica: Análisis exploratorio de datos. Fuente: notebooks/00_exploracion.ipynb, §3.1 a §3.7.
\section{Análisis exploratorio}\label{sec:exploracion}

Con los datos de la sección~\ref{sec:datos}, esta sección define primero qué se predice y con qué
métrica se mide. Después se queda con los hallazgos del notebook 00 que cambiaron una decisión; el
resto del análisis está en el notebook.

\subsection{La variable objetivo}\label{sec:objetivo}\label{sec:objetivo-formula}

Para cada reseña $i$,
\[
  Y_i =
  \begin{cases}
    1 & \text{si } \texttt{playtime\_at\_review}_i < 120 \text{ y } \texttt{voted\_up}_i = 0,\\
    0 & \text{en otro caso,}
  \end{cases}
\]
donde \texttt{playtime\_at\_review} son los minutos que llevaba jugados quien escribió la reseña
cuando la publicó y \texttt{voted\_up} es el pulgar: 1 arriba, 0 abajo\footnote{Al reseñar, Steam
pide elegir entre «Recomendado» y «No recomendado».}. En data-v1 hay \cifraPositivosEntrenamiento{}
reseñas con $Y = 1$ de \cifraResenasEntrenamiento{}, una prevalencia\footnote{La proporción de
casos positivos; aquí, de reseñas con $Y = 1$.} de \cifraPrevalenciaEntrenamientoPorResena{} por
reseña. A esa etiqueta se le llama \emph{señal de arrepentimiento temprano}. La
tabla~\ref{tab:descriptiva} resume los dos cortes de datos.

\begin{table}[htbp]
  \centering\small
  % Generado por documento/generar_figuras.py. No se edita a mano.
\begin{tabularx}{\textwidth}{>{\raggedright\arraybackslash}X rr}
\toprule
 & Entrenamiento (data-v1) & Catálogo (data-v3) \\
\midrule
Juegos & 83 & 123 \\
Reseñas & 123,972 & 184,367 \\
Reseñas con la señal ($Y = 1$) & 2,714 & 4,126 \\
Prevalencia por reseña & 2.19\,\% & 2.24\,\% \\
Reseñas publicadas & 8 abr 2023 a 14 sep 2026 & 8 abr 2023 a 21 sep 2026 \\
Publicadas desde julio de 2025 & 96.1\,\% & 97.1\,\% \\
Reseñas con perfil público & 40.1\,\% & 40.2\,\% \\
Juegos gratis & 2 & 7 \\
Precio de los de pago en pesos: mediana (mín.\ a máx.) & 400.00 (59.00 a 1,599.00) & 349.88 (56.99 a 1,599.00) \\
Nota de Metacritic: mediana (mín.\ a máx.) & 88 (61 a 97) & 86 (61 a 97) \\
Juegos sin nota de Metacritic & 23 & 33 \\
\bottomrule
\end{tabularx}

  \caption{Estadísticas descriptivas del corte de entrenamiento y del catálogo que sirve el sitio.
  El precio y la nota son de los juegos, no de las reseñas. Fuente: releases data-v1 y data-v3.}
  \label{tab:descriptiva}
\end{table}

\paragraph{Por qué es una señal proxy.}\label{sec:objetivo-proxy} Steam no registra
arrepentimiento: registra que alguien votó en contra y cuánto había jugado al escribir, así que la
señal es una proxy\footnote{Una medida que sí se puede observar y que se usa en lugar de otra que
no se puede observar directamente.}. Los 120 minutos vienen de la política de reembolsos de la
tienda (sección~\ref{sec:planteamiento}). Una reseña negativa escrita dentro de esa
ventana es lo más cercano que dejan los datos a «esto no era para mí, y me di cuenta pronto». Por
eso el sitio y este documento la llaman señal y nunca dicen que alguien se arrepintió. La señal deja
fuera a quien se arrepiente y no escribe, a quien pide el reembolso sin reseñar y a quien reseña
después de las dos horas.

\paragraph{El umbral de 120 minutos.}\label{sec:objetivo-umbral} El umbral no se eligió con los
datos, y moverlo casi no cambia qué juegos tienen más señal: con 60, 90 o 180 minutos, la
correlación de Spearman\footnote{Compara dos órdenes, no dos valores: vale 1 cuando los juegos
quedan en el mismo orden, aunque sus tasas cambien.} del orden de los juegos con el de 120 va de
\cifraSpearmanUmbralMin{} a \cifraSpearmanUmbralMax{} (detalle en el notebook 00, §3.3). Lo que sí
muestran los datos es que las negativas llegan antes: el \cifraNegativasAntesDelUmbralPorResena{}
de las reseñas negativas se escribe antes de los 120 minutos, contra el
\cifraPositivasAntesDelUmbralPorResena{} de las positivas (figura~\ref{fig:minutos}, en escala
logarítmica\footnote{Cada marca del eje multiplica a la anterior por la misma cantidad.}). En los
120 minutos no hay escalón: por minuto, las negativas pasan de
\cifraNegativasPorMinutoAntesDelUmbral{} a \cifraNegativasPorMinutoDespuesDelUmbral{} al cruzarlos,
y las positivas, de \cifraPositivasPorMinutoAntesDelUmbral{} a
\cifraPositivasPorMinutoDespuesDelUmbral{}. El único salto, a los 180 minutos y solo en las
positivas, no toca la señal.

\begin{figure}[htbp]
  \centering
  \includegraphics[width=\textwidth]{minutos-al-resenar.pdf}
  \caption{Minutos jugados al escribir la reseña, por pulgar, en escala logarítmica. Cada curva se
  normaliza sobre su propio grupo; los intervalos tienen un borde en 180 y ninguno en 120. Fuente:
  data-v1.}
  \label{fig:minutos}
\end{figure}

\subsection{Por qué PR-AUC}\label{sec:objetivo-metrica}

Con una prevalencia de \cifraPrevalenciaEntrenamientoPorResena{}, un modelo que siempre dice «no»
acierta el \cifraExactitudTrivialPorResena{} de las veces y no encuentra un solo caso, así que la
exactitud no sirve. El proyecto usa el área bajo la curva de precisión y exhaustividad
(PR-AUC)\footnote{La precisión es la parte de lo que el modelo pone arriba que sí tiene la señal; la
exhaustividad, la parte de los casos con señal que encuentra. El PR-AUC es el área bajo la curva
que traza una contra la otra.}, que con clases desbalanceadas informa más que la curva
ROC\footnote{Compara aciertos contra falsas alarmas. Con muchísimos negativos, las falsas alarmas
pesan poco y la curva puede verse bien aunque el modelo señale muchos casos de más.}
\parencite{saito2015,davis2006}. El PR-AUC de un clasificador que no sabe nada, que da a todos los
casos la misma calificación, es la prevalencia: \cifraPRAUCTrivial{}. Ese es el piso, y por eso el
PR-AUC del modelo siempre va junto al del trivial. En la validación agrupada por
juego\footnote{La validación cruzada parte los datos en partes, los \emph{folds}: entrena con todas
menos una y evalúa con la que quedó fuera, una vez por cada parte. GroupKFold reparte juegos
completos entre folds, así que mide al modelo con juegos que no vio.} el modelo llega a \cifraPRAUCModelo{}, \cifraPRAUCCociente{} veces
ese piso (sección~\ref{sec:evaluacion}). Esos cocientes contra el trivial llevan un decimal, porque
el PR-AUC cambia bastante de un fold a otro y una centésima más no diría nada; los extremos de un
intervalo llevan dos, porque ahí importa qué tan cerca quedan del 1, el valor en el que el modelo no
ordena mejor que el trivial. Todo PR-AUC lleva tres cifras significativas, y su desviación, su
intervalo o una diferencia, los mismos decimales que él. Los scores y los umbrales de las bandas
llevan cuatro decimales, porque los juegos al filo de un corte se distinguen justo en el cuarto
(sección~\ref{sec:evaluacion-filo}).

\subsection{La tasa varía por juego}\label{sec:exploracion-juego}

¿La señal es parecida en todos los juegos, o depende del juego? La
figura~\ref{fig:tasa-por-juego} ordena los \cifraJuegosEntrenamiento{} juegos de data-v1 según su
tasa de señal, la proporción de sus reseñas con $Y = 1$, con su intervalo de Wilson\footnote{El
intervalo de \textcite{wilson1927} nunca baja de cero, y por eso sirve con tasas bajas.}. La tasa va de
\cifraTasaJuegoMinPorResena{} a \cifraTasaJuegoMaxPorResena{}, y como hasta el juego con menos
reseñas tiene \cifraResenasJuegoMin{}, los intervalos son cortos al lado de esa distancia. No es
ruido: la varianza de la tasa entre juegos es unas \cifraSobredispersion{} veces la esperada si
todos tuvieran la misma tasa\footnote{Con una sola tasa para todos, la de cada juego variaría solo
por las reseñas que le tocaron, como en una binomial. El cociente divide la desviación al cuadrado
de cada juego entre esa varianza esperada y promedia: es la ji cuadrada de Pearson entre sus grados
de libertad, y vale cerca de 1 si solo el muestreo separa a los juegos.}.

\begin{figure}[htbp]
  \centering
  \includegraphics[width=\textwidth]{tasa-por-juego.pdf}
  \caption{Tasa de señal de cada juego de data-v1, de menor a mayor, con su intervalo de Wilson al
  \cifraNivelDeConfianza{}. La línea punteada es la prevalencia por reseña. Fuente: data-v1.}
  \label{fig:tasa-por-juego}
\end{figure}

Además, las cinco variables del modelo valen lo mismo en todas las reseñas de un juego: el modelo
solo puede aprender qué juegos tienen más señal. De ahí salen dos decisiones. La validación agrupa
por juego, porque con las reseñas de un juego a la vez en entrenamiento y en prueba el modelo se
mediría reconociendo juegos que ya vio (sección~\ref{sec:modelacion-validacion}). Y las variables
del juego se analizan con un punto por juego, porque \cifraResenasEntrenamiento{} filas que repiten
los mismos valores darían una precisión que no existe.

\subsection{La crítica acompaña a la señal}\label{sec:exploracion-critica}

¿Qué tienen los juegos con más señal? Juego por juego, solo la crítica tiene una correlación de
Spearman con la tasa que difícilmente sale del azar\footnote{El valor p es la probabilidad de ver
una correlación así de fuerte si en realidad no hubiera ninguna.}: \cifraSpearmanTieneNota{} con
tener nota y \cifraSpearmanNota{} con la nota, entre los juegos que la tienen, las dos con valores p
por debajo de una milésima. Con el precio es \cifraSpearmanPrecio{}, con un valor p de
\cifraSpearmanPrecioP{}, y con el descuento tampoco se distingue del azar; la gratuidad casi no se
puede medir, con \cifraGratisEntrenamiento{} juegos gratis (detalle en el notebook 00, §3.4). En la
figura~\ref{fig:tasa-contra-nota}, los \cifraSinNotaEntrenamiento{} juegos sin nota tienen una tasa
mediana de \cifraTasaSinNotaMedianaJuegos{}, contra \cifraTasaConNotaMedianaJuegos{} de los que sí
la tienen, y entre estos la nube baja conforme sube la nota, con mucha dispersión.

\begin{figure}[htbp]
  \centering
  \includegraphics[width=\textwidth]{tasa-contra-nota.pdf}
  \caption{Tasa de señal de cada juego de data-v1 contra su nota de Metacritic, en escala
  logarítmica. Los juegos sin nota van en un panel aparte, desplazados al azar a lo ancho (con
  una semilla fija); la línea negra es su mediana. Fuente: data-v1.}
  \label{fig:tasa-contra-nota}
\end{figure}

Una variable a la vez puede engañar: si los juegos caros fueran también los sin nota, el precio
parecería importar por culpa de la crítica. Con las cinco juntas, los coeficientes del modelo y un
intervalo por bootstrap sobre juegos\footnote{El bootstrap repite el ajuste con muestras de juegos
sacadas al azar y con reemplazo; el intervalo reúne el \cifraNivelDeConfianza{} central de los
coeficientes. Las variables están estandarizadas, así que los coeficientes se comparan entre sí.},
solo los intervalos de la crítica no cruzan el cero, y los dos coeficientes son negativos: tener
nota y tener una nota más alta bajan el riesgo estimado del título. El de tener nota va de
\cifraCoefTieneNotaICInf{} a \cifraCoefTieneNotaICSup{}, y el de la nota, de \cifraCoefNotaICInf{}
a \cifraCoefNotaICSup{}. El modelo se queda con las cinco variables, que se conocen antes de comprar
y valen igual para cualquier persona; lo que cambió fue la explicación, porque cada factor del sitio
lleva su nivel de evidencia, sólida para la crítica y débil para el resto
(sección~\ref{sec:interpretabilidad}).

\subsection{El texto de las negativas tempranas}\label{sec:exploracion-texto}

En el conjunto limpio (sección~\ref{sec:procesamiento-limpieza}) hay \cifraPositivasLimpias{}
positivas, \cifraNegativasTardiasLimpias{} negativas tardías y \cifraNegativasTempranasLimpias{}
negativas tempranas, las que tienen $Y = 1$. Según el log-odds con prior informativo
\parencite{monroe2008}\footnote{Compara qué tan seguido aparece cada palabra en un grupo y en el
otro; el prior evita que una palabra rara parezca distintiva solo por aparecer un par de veces.},
la palabra que más separa a las negativas tempranas de las positivas es «not», y por eso la
normalización conserva las negaciones. Frente a las tardías aparece lo que pasa al empezar:
«tutorial», «account», «login», «settings» y «minutes»; las tardías hablan de «hours», «boss»,
«mods» y «dlc». El \cifraRefundTempranasPorResena{} de las negativas tempranas menciona «refund» o
una de sus formas, contra el \cifraRefundTardiasPorResena{} de las tardías y el
\cifraRefundPositivasPorResena{} de las positivas: encaja con la lectura de la señal como proxy,
aunque no prueba que esas personas pidieran el reembolso, y la mayoría no lo menciona. El análisis
es descriptivo: el texto no entra al modelo de riesgo, y en el sitio solo lo lee la explicación de motivos,
una heurística de palabras clave (sección~\ref{sec:interpretabilidad-motivos}). La Parte A prueba
si ese lenguaje distingue de verdad a las tempranas (sección~\ref{sec:modelacion-texto}).

\subsection{Los \cifraJuegosEntrenamiento{} contra los \cifraJuegosExternos{}}\label{sec:exploracion-externos}

Los \cifraJuegosExternos{} juegos de la prueba externa pasan las mismas validaciones, pero la mezcla
es otra: más juegos gratis (\cifraGratisExternos{} contra \cifraGratisEntrenamiento{}), precios más
bajos (una mediana de \cifraPrecioMedianoExternos{} pesos contra \cifraPrecioMedianoEntrenamiento{})
y una tasa mediana por juego más alta (\cifraTasaExternosMedianaJuegos{} contra
\cifraTasaEntrenamientoMedianaJuegos{}), aunque el promedio por juego queda igual,
\cifraTasaEntrenamientoPromJuegos{} y \cifraTasaExternosPromJuegos{}
(detalle en el notebook 00, §3.6). No se decidió nada con esto: los externos nunca entran a elegir
variables, a entrenar ni a fijar umbrales. Sí cambia cómo se lee la prueba externa, que mide al
modelo con una mezcla de juegos distinta (sección~\ref{sec:evaluacion-externa}). Lo que cambió una
decisión deja dos tareas para la sección~\ref{sec:procesamiento}: limpiar el texto para el
análisis de las quejas y fijar qué variables puede usar el modelo.
